\begin{table}
    \centering
    \caption{Arhitekturna lestvica grafovskih modelov.}
    \label{tab:arhitekturna-lestvica-gnn}
    \small
    \begin{tabularx}{\textwidth}{@{}>{\raggedright\arraybackslash}p{0.24\textwidth} >{\centering\arraybackslash}p{0.10\textwidth} >{\centering\arraybackslash}p{0.12\textwidth} >{\centering\arraybackslash}p{0.10\textwidth} >{\raggedright\arraybackslash}X@{}}
        \toprule
        \textbf{Model} & \textbf{Ločene relacije} & \textbf{Fuzija relacij} & \textbf{Uteži povezav} & \textbf{Namen stopnje} \\
        \midrule
        \ModelName{GCN} & ne & -- & ne & preveri osnovno grafovsko predstavitev v homogenem grafu \\
        \ModelName{HeteroGCN-mean} & da & povprečje & ne & preveri vpliv ločene obravnave relacij \\
        \ModelName{HeteroGCN-MLP} & da & \ModelName{MLP} & ne & preveri vpliv naučene fuzije relacijskih predstavitev \\
        \ModelName{HeteroGCN-MLP-w} & da & \ModelName{MLP} & da & preveri vpliv naučenih predznačenih uteži povezav \\
        \bottomrule
    \end{tabularx}
\end{table}
